\scanpage{9}
\mspage{8}
\deleted{
\uline{Question} Comment définir directement une connexion intégrable sur $\mathbb R^*f_*(\Omega^\bullet_{X/S})$ ??

\uline{Lemme}\quad $X$ Topos, $K^\bullet,L^\bullet$ complexes abéliens de $X$, $\unsure{\mathcal G|\mathfrak G}$ faisceau abélien de $X$,
\diagram{\begin{tikzcd}[ampersand replacement=\&]
\unsure{\mathcal G|\mathfrak G}\arrow[r,"i"]\&\underline{\operatorname{Hom}}(K^\bullet,L^\bullet)^{-1}=\prod_i\underline{\operatorname{Hom}}(K^i,L^{i-1})
\end{tikzcd}}
$P$ un $\unsure{\mathcal G|\mathfrak G}$-torseur,
\diagram{\begin{tikzcd}[ampersand replacement=\&]P\arrow[r,"\Theta"]\&\mathrm Z(\underline{\operatorname{Hom}}(K^\bullet,L^\bullet)^{(0)})\end{tikzcd}}
\deleted{Alors on peut associer à $P$ un homomorphisme de $D(X)$, $\unsure{\alpha\circ\varphi}$}
\diagram{\begin{tikzcd}[ampersand replacement=\&]\deleted{K^\bullet}\arrow[r]\&\deleted{\unsure{L^\bullet}}\end{tikzcd}}
un homomorphisme compatible avec les opérations de $\underline{\operatorname{Hom}}(K^\bullet,L^\bullet)^{-1}$
\[
\Theta(g.p)=\Theta(p)+\underbrace{(d\circ i(g)+i(g)\circ d)}_{\delta(i(g))}
\]
Alors on peut trouver un hom \uline{dans} $D(X)$
\diagram{\begin{tikzcd}[ampersand replacement=\&]
\varphi=\varphi(\inserted{K,L,}\unsure{\mathcal G|\mathfrak G},P,i,\Theta):\quad K^\bullet\arrow[r]\&L^\bullet
\end{tikzcd}}
tel que :

1°) $\varphi$ est fonctoriel en $L^\bullet$ (en un sens évident)
\margin{\deleted{complexe} variable dans $C^+(X)$}

2°) \deleted{Si $P$ est \unsure{trivial}\ \unsure{trivial évident}} le $=$ si on remplace $\unsure{\mathcal G|\mathfrak G},P,i,\Theta$ par $\unsure{\mathcal G|\mathfrak G}',P',i',\Theta'$, avec
\diagram{\begin{tikzcd}[ampersand replacement=\&]\unsure{\mathcal G|\mathfrak G}\arrow[r,"u"]\&\unsure{\mathcal G|\mathfrak G}'\&P\arrow[r,"v"]\&P'\end{tikzcd}}
pour $i=i'\circ u$, $\Theta=v\Theta'$, $v$ \deleted{\unsure{compatible aux}} \unsure{un $u$-hom.}

3°) Si $P$ admet une section, alors
}
